\documentclass[aps,prd,onecolumn,longbibliography,nofootinbib]{revtex4-2}

\usepackage[utf8]{inputenc}
\usepackage{amsmath,amssymb,mathtools}
\usepackage{graphicx}
\usepackage{xcolor}
\usepackage{hyperref}
\usepackage{enumitem}
\usepackage{booktabs}
\usepackage[final]{microtype}

\hypersetup{colorlinks=true, linkcolor=blue, citecolor=blue, urlcolor=blue}

\newcommand{\dd}{\mathrm{d}}
\newcommand{\Tr}{\mathrm{Tr}}
\newcommand{\Htot}{\mathcal{H}_{\rm tot}}
\newcommand{\Hcore}{\mathcal{H}_{\rm core}}
\newcommand{\Hrad}{\mathcal{H}_{\rm rad}}
\newcommand{\Svn}{S_{\rm vN}}
\newcommand{\Smicro}{S_{\rm micro}}
\newcommand{\Rn}{\mathcal{R}}

\begin{document}

\title{An Abstract Sector Hamiltonian for Black-Hole-Like Evaporation Phenomenology}

\author{[Author]}
\affiliation{[Affiliation]}
\date{\today}

\begin{abstract}
We describe a finite non-gravitational quantum model that reproduces a large
part of the phenomenological package associated with black-hole evaporation.
The model is an abstract sector Hamiltonian.  The core Hilbert space is a
direct sum of finite sectors with an area-like state count, the mass assigned
to sector size obeys a Schwarzschild-like relation, and radiation is produced
by Hamiltonian couplings from one core sector to a smaller core sector plus
outgoing modes.  Golden-rule transition rates are then determined by matrix
elements, radiation phase space, and the density of final core states.  In an
energy-resolved implementation, an exponential local density of states gives a
near-thermal hard spectrum, while within-sector mixing restores accelerating
evaporation for the black-hole-like mass law; a linear mass-law control does
not show the same acceleration.  A companion repeated-interaction realization
tracks purification and shows, in small exact state-vector runs, a shrinking
core capacity, soft information transfer, enhanced early/late radiation mutual
information with scrambling, and a finite Page-like turnover of the full
radiation von Neumann entropy.  The construction does not derive the area
state count, the mass law, or the mixing dynamics from microscopic gravity.
Its purpose is narrower: to give a reproducible non-gravitational quantum
comparison model in which thermality, negative heat capacity, accelerating
emission, and unitary information flow can be assembled and separately tested.
\end{abstract}

\maketitle

\section{Introduction}

Black-hole evaporation combines several statements that are often discussed
together.  A Schwarzschild black hole has entropy proportional to area, a
temperature inversely proportional to mass, negative heat capacity, and an
evaporation rate that increases as the mass decreases.  Hawking's calculation
also gives locally thermal radiation from a quantum field on a black-hole
background \cite{Hawking1975}.  Page's argument gives the expected
fine-grained radiation entropy for unitary evaporation
\cite{Page1993Info,Page1993Average}.  Modern gravitational calculations
recover this fine-grained behavior using generalized entropy, quantum extremal
surfaces, islands, and replica-wormhole saddles
\cite{Penington2020,AEMM2019,AlmheiriReview2021}.

This paper asks a narrower question.  How much of the usual evaporation
phenomenology can be reproduced by an ordinary finite quantum model with no
spacetime dynamics?  The target is a comparison model.  It should be explicit
enough that its Hilbert space, evolution rule, initial state class, observables,
and controls are all inspectable.  The model should also keep separate the
different entropies involved: the microcanonical entropy of the core, the
entropy of the hard radiation marginal, and the von Neumann entropy of the full
emitted radiation.

Toy models of black-hole evaporation have a long history.  Qubit and circuit
models isolate unitary information transfer and Page-like behavior
\cite{Avery2013,OsugaPage2018,Broda2021}.  Tensor-network and random-unitary
models define explicit maps from black-hole microstates to black-hole plus
radiation Hilbert spaces \cite{LeutheusserVanRaamsdonk2017,
PiroliSunderhaufQi2020}.  The model studied here is closest in spirit to those
algorithmic constructions.  Its central emphasis is thermodynamic: the same
finite model is required to carry a black-hole-like entropy/mass relation,
negative heat capacity, a generated emission rate, and unitary information
flow.

The most compact version is an abstract sector Hamiltonian.  The core Hilbert
space is a direct sum
\begin{equation}
    \Hcore=\bigoplus_n B_n ,
\end{equation}
with sector dimensions chosen so that
\begin{equation}
    \Smicro(n)=\log\dim B_n
\end{equation}
has the desired area-like scaling.  Radiation modes are coupled to transitions
from $B_n$ to $B_{n-1}$.  The transition rates are generated by Hamiltonian
matrix elements and final-state phase space.  The model is finite and
non-gravitational; the area-like sector count and mass law are inputs.

The main result is constructive.  In a sector Hamiltonian with
energy-resolved core sectors, the combination of a black-hole-like mass law,
an exponential local density of states, and within-sector mixing gives both a
near-thermal hard spectrum and accelerating evaporation.  The same construction
with a linear mass law gives a near-thermal hard spectrum but does not show the
same acceleration.  A separate exact repeated-interaction realization adds
purification registers and soft records, demonstrating the compatible
information-flow structure at small finite size.

\section{Model Standard and Entropies}
\label{sec:standard}

The comparison standard is:
\begin{enumerate}[leftmargin=*]
    \item specify a finite Hilbert space;
    \item specify a Hamiltonian or Hamiltonian-derived transition rule;
    \item specify an initial state class;
    \item define the core, radiation, bath/purifier, and record subsystems;
    \item measure thermodynamic and information-flow observables.
\end{enumerate}

Three entropies are used.  The microcanonical or Boltzmann entropy of the core
is
\begin{equation}
    \Smicro(n)=\log\Omega_n,\qquad \Omega_n=\dim B_n .
\end{equation}
This entropy is used to define the thermodynamic temperature and heat
capacity.  The hard-radiation entropy is the entropy of the visible hard
radiation marginal after hidden bath or purifier records are traced out.  The
Page-like entropy is the von Neumann entropy of the full emitted radiation
subsystem,
\begin{equation}
    \Svn(R)=-\Tr \rho_R\log\rho_R .
\end{equation}
The second Renyi entropy $S_2(R)=-\log\Tr\rho_R^2$ is useful for larger
numerical diagnostics, but the Page-like results reported below use von
Neumann entropy.

\section{Abstract Sector Hamiltonian}
\label{sec:hamiltonian}

The total Hilbert space is
\begin{equation}
    \Htot=\left(\bigoplus_{n=n_{\min}}^{n_{\max}} B_n\right)\otimes \Hrad ,
\end{equation}
where $B_n$ is the active core sector and $\Hrad$ is a finite radiation
Fock-space truncation or a finite set of outgoing wave-packet modes.  A minimal
Hamiltonian is
\begin{equation}
    H_{\rm tot}=H_{\rm core}+H_{\rm rad}+H_{\rm int}.
\end{equation}
The core Hamiltonian has block form
\begin{equation}
    H_{\rm core}
    =
    \sum_{n,a} E_{n,a}\,
    |n,a\rangle\langle n,a| ,
\end{equation}
with $a=1,\ldots,\dim B_n$.  We write
\begin{equation}
    E_{n,a}=M_n+\epsilon_{n,a}.
\end{equation}
The sector label controls the coarse mass scale $M_n$, while
$\epsilon_{n,a}$ is an internal excitation energy.

The radiation Hamiltonian is
\begin{equation}
    H_{\rm rad}=\sum_\lambda \omega_\lambda
    b_\lambda^\dagger b_\lambda .
\end{equation}
The interaction Hamiltonian couples a core state in sector $n$ to a core state
in sector $n-1$ plus one outgoing radiation quantum:
\begin{equation}
    H_{\rm int}
    =
    \sum_{n,a,b,\lambda}
    \left(
        g^{(\lambda)}_{n;b a}
        |n-1,b\rangle\langle n,a|\, b_\lambda^\dagger
        +{\rm h.c.}
    \right).
\end{equation}
In the weak-coupling or secular limit, the transition rate is
\begin{equation}
    \Gamma_{n,a\rightarrow n-1,b;\lambda}
    =
    2\pi
    \left|g^{(\lambda)}_{n;b a}\right|^2
    \delta(E_{n,a}-E_{n-1,b}-\omega_\lambda).
\end{equation}
With a continuum or binned radiation sector this becomes the usual
golden-rule expression: matrix element times final-state density.

This Hamiltonian is abstract.  Its role is to replace an externally assigned
evaporation schedule by Hamiltonian sector transitions whose rates can be
measured and controlled.

\section{Thermodynamic Dictionary}
\label{sec:thermo}

Two equivalent labels are useful.  Let $n$ be an area-like label and $L$ be a
linear size label with $n=L^2$.  The sector count is
\begin{equation}
    \dim B_n=q^n,\qquad \Smicro(n)=n\log q .
\end{equation}
The black-hole-like mass law is
\begin{equation}
    M_n=\alpha\sqrt n .
\end{equation}
Then
\begin{equation}
    \frac{\dd \Smicro}{\dd M}
    =
    \frac{\dd \Smicro/\dd n}{\dd M/\dd n}
    =
    \frac{\log q}{\alpha/(2\sqrt n)}
    =
    \frac{2\sqrt n\log q}{\alpha},
\end{equation}
and
\begin{equation}
    T_n=\left(\frac{\dd \Smicro}{\dd M}\right)^{-1}
    =
    \frac{\alpha}{2\sqrt n\log q}
    \sim M^{-1}.
\end{equation}
The heat capacity is negative because $M$ decreases as $T$ increases.

The linear control is
\begin{equation}
    M_n=\alpha n ,
\end{equation}
which gives a constant temperature for the same sector count.  This control is
important because it keeps the finite-sector framework while removing the
black-hole-like mass/entropy relation.

For an emitted quantum of energy $\omega$, the microcanonical final-state
factor is
\begin{equation}
    \exp[\Smicro(M-\omega)-\Smicro(M)] .
\end{equation}
For $\omega\ll M$ this gives
\begin{equation}
    \Smicro(M-\omega)-\Smicro(M)
    =
    -\beta(M)\omega+O(\omega^2),
\end{equation}
with $\beta=1/T$.  A hard radiation spectrum therefore becomes thermal when
the core density of states and transition matrix elements are sufficiently
smooth across the relevant energy window.

\section{Energy-Resolved Sectors}
\label{sec:energy-resolved}

The first sector-Hamiltonian diagnostic used small random perturbations inside
each sector and generated accelerating evaporation for the square-root mass
law.  The hard spectrum was too narrow.  The improved construction makes
the internal density of states explicit.  Each sector has energies
\begin{equation}
    E_{n,a}=M_n+\epsilon_{n,a},
\end{equation}
where the local density of states is approximately exponential in an energy
window of width $O(T_n)$:
\begin{equation}
    \rho_n(\epsilon)\propto e^{\beta_n \epsilon}.
\end{equation}

Transitions are generated by shrinkage operators
\begin{equation}
    X_n:B_n\rightarrow B_{n-1},
\end{equation}
with rates
\begin{equation}
    \Gamma_{f i}\propto
    |\langle f,n-1|X_n|i,n\rangle|^2\,
    \omega_{fi}^{p},
    \qquad
    \omega_{fi}=E_i^{(n)}-E_f^{(n-1)}.
\end{equation}
The exponent $p$ represents the radiation phase-space or ohmic weight used in
the diagnostic.  The emitted spectrum is compared in the variable
\begin{equation}
    x=\beta_n\omega
\end{equation}
to the target number distribution
\begin{equation}
    P(x)\propto x^p e^{-x}.
\end{equation}

The population evolution used in the numerical diagnostic is the secular
rate equation
\begin{equation}
    \dot p_{n,i}
    =
    -\sum_{f}\Gamma_{fi}p_{n,i}
    +
    \sum_{j}\Gamma_{ij}p_{n+1,j}.
\end{equation}
This is the Markovian rate limit of the Hamiltonian above.  The fully coherent
state-vector realization with explicit radiation wave packets is a separate
upgrade.

\section{Numerical Results}
\label{sec:results}

\subsection{Coarse sector Hamiltonian}

The coarse sector-Hamiltonian diagnostic used
\begin{equation}
    n=4,\ldots,10,\qquad q=2,\qquad \alpha=8,
\end{equation}
with local and scrambled shrinkage maps, two seeds, and a finite emission
passband.  For the square-root mass law, emitted power accelerated.  For the
linear mass law, emitted power decelerated.  The local and scrambled removal
maps gave nearly identical coarse rates, showing that the mass/entropy sector
structure dominated that diagnostic.

The main weakness of the coarse run was the hard spectrum.  The total variation
distance from the target thermal distribution in $x=\beta\omega$ was about
$0.77$ for the square-root mass law.  This motivated the energy-resolved sector
construction.

\subsection{Energy-resolved sector Hamiltonian}

With exponential sector density of states and no within-sector
rethermalization, the hard spectrum became close to thermal.  The power
decelerated because the state population drifted into lower-emission parts of
the sector spectrum.

The next diagnostic applied complete within-sector mixing after each emission
step:
\begin{equation}
    p_n(a)\mapsto \frac{P_n}{\dim B_n}.
\end{equation}
This is an idealized stand-in for fast intra-sector scrambling.  With this
mixing, the square-root mass law produced both a near-thermal hard spectrum and
accelerating evaporation.

For two seeds, local and scrambled shrinkage, and width parameter
${\rm width}_x=4$, the square-root mass law gave:
\begin{center}
\begin{tabular}{lccc}
\toprule
operator & power ratio & jump ratio & TV distance\\
\midrule
local     & 1.098 & 1.045 & 0.055\\
scrambled & 1.097 & 1.045 & 0.054\\
\bottomrule
\end{tabular}
\end{center}
Here the power ratio compares mid-evaporation power to early power.  Values
above one indicate acceleration.  The total variation distance compares the
emitted number distribution to $x^p e^{-x}$.

The corresponding linear mass-law control had a similarly good spectrum but
did not accelerate:
\begin{center}
\begin{tabular}{lccc}
\toprule
operator & power ratio & jump ratio & TV distance\\
\midrule
local     & 0.988 & 0.995 & 0.051\\
scrambled & 0.988 & 0.995 & 0.049\\
\bottomrule
\end{tabular}
\end{center}

The interpretation is direct.  The exponential local density of states fixes
the hard spectrum.  Within-sector mixing prevents the core from falling into
low-emission states.  The square-root mass law then gives acceleration, while
the linear control remains approximately flat or mildly decelerating.

\subsection{Repeated-interaction information-flow realization}

The sector-Hamiltonian rate model does not by itself track purification.  A
companion finite repeated-interaction model does.  It uses core sectors
$B_L$ with
\begin{equation}
    \dim B_L=q^{L^2},
\end{equation}
hard radiation registers, hidden bath/purifier registers, soft shell records,
and an emitted-energy accumulator.  The exact state-vector runs are small:
\begin{equation}
    L_{\rm init}=3,\qquad q=2,\qquad n=6
\end{equation}
for the main information-flow diagnostic.

The hard marginal is fixed by the chosen thermodynamic emission weights and is
purified by the bath registers.  The soft records carry shell information when
the active core capacity shrinks.  With scrambling, the soft entropy and
old/new radiation mutual information increase relative to the no-scrambling
control.

For the default six-emission run, the full old/new radiation mutual
information was:
\begin{center}
\begin{tabular}{lc}
\toprule
scrambler & $I({\rm old}:{\rm new})$\\
\midrule
Margulis-like & 2.615\\
grid-like     & 2.624\\
none          & 1.988\\
\bottomrule
\end{tabular}
\end{center}
The enhancement is about $0.63$ nats.

At threshold $\Delta=4$, a fixed eight-emission hard schedule gives a small
Page-like turnover in the von Neumann entropy of the full emitted radiation:
\begin{center}
\begin{tabular}{ccccc}
\toprule
emissions & 2 & 4 & 6 & 8\\
\midrule
scrambled $\Svn(R_{\rm full})$ & 1.427 & 3.799 & 2.612 & 1.900\\
no-scrambling $\Svn(R_{\rm full})$ & 1.017 & 1.384 & 1.592 & 1.698\\
\bottomrule
\end{tabular}
\end{center}
This is a finite small-system diagnostic.  The thermodynamic-limit Page curve
remains a larger-size target.

\section{Controls and Mechanism Separation}
\label{sec:controls}

The current evidence separates several mechanisms.

Changing the mass law distinguishes the source of acceleration.  With the same
sector framework and a good hard spectrum, the square-root mass law accelerates
while the linear mass law does not.  This ties acceleration to the
black-hole-like entropy/mass relation.

Changing the sector density of states distinguishes spectrum generation from
rate generation.  A coarse internal spectrum gives acceleration and a poor
hard spectrum.  An exponential local density of states gives a near-thermal
hard spectrum.  Without mixing, that improved spectrum comes with deceleration;
with mixing, acceleration returns.

Changing scrambling in the repeated-interaction realization distinguishes hard
local thermality from information transfer.  The hard marginal remains fixed
by the emission channel, while soft entropy and old/new mutual information
depend on scrambling.

The strongest remaining modeled ingredient is within-sector mixing.  In the
present diagnostic it is represented by the idealized map
\begin{equation}
    p_n(a)\mapsto P_n/\dim B_n .
\end{equation}
In a more microscopic model, this should arise from chaotic intra-sector
Hamiltonian dynamics.

\section{Relation to a Natural Microscopic Model}
\label{sec:natural}

The abstract sector Hamiltonian is a diagnostic construction rather than a
natural microscopic derivation.  It specifies sectors, state counts, mass laws,
and transition operators.  The recent connector-model search gives one possible
direction for reducing those inputs.  A core with $N$ active constituents and
$O(N^2)$ relational connector modes naturally gives an area-like state count.
The viable spectral target is now narrower: a protected quadratic
($z=2$) connector band, reminiscent of type-B Goldstone or ferromagnetic
magnon physics, plus an emission coupling with
\begin{equation}
    |g(\omega)|^2\sim \omega^{1/2}
\end{equation}
to recover the desired $P\sim M^{-2}$ scaling for that band.

This connector route remains open.  It would need to explain the ordering of
$O(N^2)$ connector modes, the protected quadratic dispersion, the emission
matrix-element scaling, and slow shrinkage after many microscopic emissions.
Until those ingredients arise from one recognizable Hamiltonian, the sector
Hamiltonian is best viewed as an existence and diagnostic model rather than a
natural microscopic evaporator.

\section{Discussion}

The construction gives a finite non-gravitational comparison model for a
substantial part of black-hole evaporation phenomenology.  The thermodynamic
sector structure supplies $S\sim M^2$, $T\sim M^{-1}$, and negative heat
capacity.  Hamiltonian-derived sector transitions generate rates.  An
energy-resolved sector density of states and within-sector mixing produce a
near-thermal hard spectrum together with accelerating evaporation for the
black-hole-like mass law.  A companion exact repeated-interaction model shows
that this thermodynamic behavior is compatible with unitary purification,
scrambling-dependent soft information transfer, and a small Page-like
radiation entropy turnover.

The limitations are essential to the claim.  The area-like sector dimensions
are specified.  The black-hole-like mass law is specified.  The current
energy-resolved diagnostic uses idealized within-sector mixing.  The exact
information-flow runs are small.  The fully coherent autonomous Hamiltonian
with outgoing radiation wave packets remains to be implemented.

The useful conclusion is that the evaporation package can be decomposed into
finite quantum statistical mechanisms.  The sector model shows which parts can
be assembled without dynamical spacetime, and which parts remain tied to model
inputs.  A natural microscopic derivation of the sector count, the mass law,
and the mixing dynamics would be a stronger result.

\appendix

\section{Reproducibility Notes}

The main sector-Hamiltonian diagnostics are:
\begin{center}
\begin{tabular}{ll}
\toprule
script & result note\\
\midrule
\texttt{sim/sector\_hamiltonian\_evaporator.py} &
\texttt{notes/sector\_hamiltonian\_evaporator\_results.md}\\
\texttt{sim/sector\_hamiltonian\_energy\_resolved.py} &
\texttt{notes/sector\_hamiltonian\_energy\_resolved\_results.md}\\
\texttt{sim/fused\_floquet\_time\_resolved\_scan.py} &
\texttt{notes/fused\_floquet\_time\_resolved\_results.md}\\
\texttt{sim/repeated\_interaction\_page\_trajectory.py} &
\texttt{sim/data/repeated\_interaction\_page\_trajectory.csv}\\
\bottomrule
\end{tabular}
\end{center}

The energy-resolved sector-Hamiltonian runs use:
\begin{equation}
    n=4,\ldots,10,\qquad q=2,\qquad \alpha=8,
\end{equation}
two seeds, local and scrambled shrinkage operators, exponential sector density
of states, and width parameters ${\rm width}_x=2,4,8$.  The reported tables use
${\rm width}_x=4$ unless stated otherwise.

The repeated-interaction information-flow runs use:
\begin{equation}
    L_{\rm init}=3,\qquad q=2,
\end{equation}
six emissions for the main old/new mutual information diagnostic and eight
emissions for the Page-like trajectory.  The old radiation consists of the
first half of the emitted records and the new radiation of the second half.
The full radiation subsystem includes hard, bath/purifier, and soft records.

To regenerate the central CSV outputs from the project root, run:
\begin{verbatim}
python bh-evaporator/sim/sector_hamiltonian_evaporator.py
python bh-evaporator/sim/sector_hamiltonian_energy_resolved.py
python bh-evaporator/sim/fused_floquet_time_resolved_scan.py
python bh-evaporator/sim/repeated_interaction_page_trajectory.py
\end{verbatim}
The scripts use standard Python scientific packages, primarily
\texttt{numpy} and \texttt{scipy}.  Once the stated seeds are fixed, the
diagnostics are deterministic.

\bibliographystyle{apsrev4-2}
\bibliography{refs}

\end{document}
